\section{Dynamic-Solver Arguments}

This appendix supplies the dynamic-solver details used in
Section~\ref{sec:model}. It is organized around the three properties needed for
the categorical reduction. First, locally well-posed projected trajectories give
a terminal map. Second, that terminal map is measurable. Third, ordinary
terminal outcomes lie among the roots of the projected vector field.
The first subsection proves the sufficient condition used in
Lemma~\ref{lemma:Lipschitz}: locally Lipschitz dynamics generate a unique
maximal local projected flow. The second subsection proves that the terminal
map $g_{\vQ}$ is measurable even though some trajectories may fail to exist for
all time or may fail to converge. The third subsection proves the root
containment behind Proposition~\ref{prop:finite-rto-roots}. This last argument
is included because $\QTC$ need not be continuous at the boundary, so standard
continuous-vector-field arguments such as Barbalat's lemma do not apply
directly.

\subsection{Well-Posedness Under Local Lipschitz Dynamics}\label{proof:lipschitz-local-flow}

This subsection proves Lemma~\ref{lemma:Lipschitz}. The global Lipschitz case
is a standard projected-dynamical-system theorem. The locally Lipschitz case
is obtained by the global through the following trick: for any compact
$K\subseteq X$ there exists a globally Lipschitz map $\widehat{\vQ}_K$ that
agrees with $\vQ$ on $K$. As the local behavior of the flow depends only on the
local dynamics, we can use the global flow $\widehat{\vG}_K$ of
$\widehat{\vQ}_K$ to define the local flow of $\vQ$ until the exit time from
$K$. By patching together these local flows, we can construct a maximal local
flow for $\vQ$ on $X$. The uniqueness and continuity properties follow from the
uniqueness and continuity of the global flows on each compact region and the
consistency of these flows before exit, verifying
Assumption~\ref{ass:well-behaved-proj}.


\begin{lemma}[Global well-posedness under global Lipschitzness]\label{lemma:global-lip-flow}
If $\vQ$ is globally Lipschitz on $X$, then the projected dynamics generated by
$\vQ$ admit a unique global flow
\[
\vG:X\times[0,\infty)\to X
\]
that is continuous in $(\vx,t)$.
\end{lemma}
\begin{proof}
Global Lipschitzness implies the one-sided Lipschitz and linear-growth
conditions required by Theorem~2.5 of \citet{nagurney2012projected} on a closed
convex set. Their projected system uses the opposite sign convention, so after
replacing their vector field by ours the theorem applies to
$\dot \vx=\Pi_{\TC_X(\vx)}\vQ(\vx)$. It gives a unique global solution and
continuous dependence on initial conditions, uniformly on compact time
intervals. Together with continuity of each trajectory in time, this yields
continuity of the flow map $(\vx,t)\mapsto \vG(\vx,t)$. In the notation of
Assumption~\ref{ass:well-behaved-proj}, the global case corresponds to
$\tau(\vx)=\infty$ for every $\vx\in X$.
\end{proof}

We now pass from global Lipschitzness to local Lipschitzness by localization.
The argument has five steps. First, Lemma~\ref{lemma:global-lip-flow} covers the
case in which $\vQ$ is globally Lipschitz. Second, for each bounded region
$K_R$, Lemma~\ref{lemma:localized-global-systems} constructs a globally
Lipschitz auxiliary dynamics $\widehat{\vQ}_{K_R}$ that agrees with $\vQ$ on
$K_R$; this auxiliary dynamics has a global flow
$\widehat{\vG}_{K_R}$. Third, Lemma~\ref{lemma:localized-consistency} shows that
these localized flows are compatible before exit, so the construction does not
depend on arbitrary choices of extensions. Fourth,
Lemma~\ref{lemma:local-flow-construction} patches the localized flows into a
single local solution map on all of $X$. Fifth,
Lemma~\ref{lemma:maximal-local-flow} shows that the patched solution map is the
unique maximal local flow of the original projected dynamics.

\begin{lemma}[Lipschitz extensions on bounded regions]\label{lemma:localized-global-systems}
Suppose $\vQ$ is locally Lipschitz. For each $R>0$, there exists a globally
Lipschitz map $\widehat{\vQ}_{K_R}:X\to\R^d$ such that
$\widehat{\vQ}_{K_R}=\vQ$ on $K_R:=X\cap\overline B_R$($B_R$ is the open ball of radius $R$ centered at the origin). The projected dynamics
generated by $\widehat{\vQ}_{K_R}$ have a unique global flow
$\widehat{\vG}_{K_R}:X\times[0,\infty)\to X$ that is continuous in $(\vx,t)$.
\end{lemma}
\begin{proof}
Local Lipschitzness is understood relative to $X$; by compactness, a finite
subcover of local Lipschitz neighborhoods gives a single Lipschitz constant for
$\vQ|_{K_R}$. Applying the McShane extension theorem coordinate by coordinate
gives a globally Lipschitz map
\[
\widehat{\vQ}_{K_R}:X\to\R^d
\]
possibly with a larger Lipschitz constant, such that
$\widehat{\vQ}_{K_R}=\vQ$ on $K_R$. The conclusion follows from
Lemma~\ref{lemma:global-lip-flow} applied to $\widehat{\vQ}_{K_R}$.
\end{proof}

For the rest of this subsection, fix the maps $\widehat{\vQ}_{K_R}$ and flows
$\widehat{\vG}_{K_R}$ from Lemma~\ref{lemma:localized-global-systems}. Define
the exit time
\[
\sigma_{K_R}(\vx):=\inf\{t\ge 0:\|\widehat{\vG}_{K_R}(\vx,t)\|\ge R\},
\]
with the convention $\inf\varnothing=\infty$.
Thus $\sigma_{K_R}(\vx)$ is the first time the global flow generated by
$\widehat{\vQ}_{K_R}$ leaves the ball $B_R$ used to define the
localization. Before that time, the localized vector field agrees with the
original vector field.

The localized extensions are chosen separately, so their global flows need not
agree after they leave the region where they both coincide with $\vQ$. The next
lemma records the compatibility needed for patching: before a smaller
localization exits, its flow is exactly the restriction of any larger
localization.

\begin{lemma}[Agreement before exit]\label{lemma:localized-consistency}
Fix $R>0$, $\vx\in X$, and $T<\sigma_{K_R}(\vx)$. Then
$\widehat{\vG}_{K_R}(\vx,\cdot)$ solves the original projected dynamics on
$[0,T]$.

Consequently, for every $S>R$,
\[
\widehat{\vG}_{K_R}(\vx,t)=\widehat{\vG}_{K_S}(\vx,t)
\qquad \text{for all }t\in[0,T].
\]
\end{lemma}
\begin{proof}
Since $T<\sigma_{K_R}(\vx)$,
$\widehat{\vG}_{K_R}(\vx,t)\in K_R$ for every $t\in[0,T]$. Write
$\vy(t):=\widehat{\vG}_{K_R}(\vx,t)$. Hence
\[
\widehat{\vQ}_{K_R}(\vy(t))=\vQ(\vy(t))
\qquad \text{for all }t\in[0,T].
\]
Therefore $\widehat{\vG}_{K_R}(\vx,\cdot)$ solves the original projected
dynamics on $[0,T]$.

Now let $S>R$. Since $K_R\subseteq K_S$, the same trajectory remains in the
region $K_R$, where both localized vector fields agree with $\vQ$. Therefore
along $\vy(\cdot)$, both
$\widehat{\vQ}_{K_R}(\vy(t))=\vQ(\vy(t))$ and
$\widehat{\vQ}_{K_S}(\vy(t))=\vQ(\vy(t))$ hold for all $t\in[0,T]$.
Thus the path $\vy(\cdot)$ also solves the globally Lipschitz projected
dynamics generated by $\widehat{\vQ}_{K_S}$ on $[0,T]$, with initial condition
$\vy(0)=\vx$. The flow
$\widehat{\vG}_{K_S}(\vx,\cdot)$ is the unique solution of that same
$\widehat{\vQ}_{K_S}$-system from the same initial condition. Hence
\[
\vy(t)=\widehat{\vG}_{K_S}(\vx,t)
\qquad \text{for all }t\in[0,T].
\]
Since $\vy(t)=\widehat{\vG}_{K_R}(\vx,t)$, the desired agreement follows.
\end{proof}

\begin{lemma}[Constructing the maximal local flow]\label{lemma:local-flow-construction}
Define
\[
\tau(\vx):=\sup_{R>\|\vx\|}\sigma_{K_R}(\vx).
\]
If $t<\tau(\vx)$, then some $R>\|\vx\|$ satisfies
$t<\sigma_{K_R}(\vx)$; otherwise all exit times in the supremum would be at
most $t$, contradicting $t<\tau(\vx)$.
For $0\le t<\tau(\vx)$, choose $R>\|\vx\|$ such that
$t<\sigma_{K_R}(\vx)$ and set
\[
\vG(\vx,t):=\widehat{\vG}_{K_R}(\vx,t).
\]
Then $\tau(\vx)>0$, $\vG$ is well defined on
$D:=\{(\vx,t):0\le t<\tau(\vx)\}$, and $\vG(\vx,\cdot)$ solves the original
projected dynamics on every compact subinterval of $[0,\tau(\vx))$.
\end{lemma}
\begin{proof}
The lifetime is strictly positive: if $R>\|\vx\|$, continuity of
$t\mapsto \widehat{\vG}_{K_R}(\vx,t)$ at $t=0$ gives
$\sigma_{K_R}(\vx)>0$.

Lemma~\ref{lemma:localized-consistency} shows that the value of
$\vG(\vx,t)$ is independent of which such localization is chosen. To see that
$\vG(\vx,\cdot)$ solves the original dynamics on compact subintervals, fix
$T<\tau(\vx)$. By the same supremum argument, choose $R>\|\vx\|$ with
$T<\sigma_{K_R}(\vx)$. Then
$\vG(\vx,t)=\widehat{\vG}_{K_R}(\vx,t)$ for all $t\in[0,T]$ by
Lemma~\ref{lemma:localized-consistency}, and
$\widehat{\vG}_{K_R}(\vx,\cdot)$ solves the original projected dynamics on
$[0,T]$.
\end{proof}

\begin{lemma}[Maximality and uniqueness]\label{lemma:maximal-local-flow}
The lifetime $\tau$ from Lemma~\ref{lemma:local-flow-construction} is the
maximal existence time. Moreover, any local solution of the original projected
dynamics agrees with $\vG$ on its interval of existence; hence $\vG$ is the
unique maximal local solution map.
\end{lemma}
\begin{proof}
Any other local solution stays in some bounded region on each compact time
interval. On that region it is also a solution of one globally Lipschitz
localized system, where uniqueness is already known.
Let $\vy:[0,T]\to X$ be any solution of the original
projected dynamics with $\vy(0)=\vx$. Since $\vy([0,T])$ is compact, choose
$R$ such that
\[
\sup_{0\le t\le T}\|\vy(t)\|<R.
\]
Then $\widehat{\vQ}_{K_R}=\vQ$ along $\vy([0,T])$, so $\vy$ also solves the
globally Lipschitz projected dynamics generated by $\widehat{\vQ}_{K_R}$. By
uniqueness for the $\widehat{\vQ}_{K_R}$-system,
\[
\vy(t)=\widehat{\vG}_{K_R}(\vx,t)\qquad\text{for all }t\in[0,T].
\]
Because this path agrees with $\vy$ and stays strictly inside $B_R$ on $[0,T]$,
we have $T<\sigma_{K_R}(\vx)$, and therefore $T<\tau(\vx)$. This proves
uniqueness on each compact interval on which a solution exists. Applying the
same argument on each compact interval $[0,T]\subset[0,S)$ shows that every
solution on $[0,S)$ satisfies $S\le \tau(\vx)$. Since
Lemma~\ref{lemma:local-flow-construction} already constructs a solution on
$[0,\tau(\vx))$, $\tau(\vx)$ is the maximal existence time.
\end{proof}

\begin{lemma}[Regularity of $\tau$ and $\vG$]\label{lemma:lsc-flow-continuity}
The maximal existence time $\tau$ from Lemma~\ref{lemma:local-flow-construction} is
lower semicontinuous, and the local flow $\vG$ is continuous on
$D:=\{(\vx,t):0\le t<\tau(\vx)\}$.
\end{lemma}
\begin{proof}
We first prove lower semicontinuity of $\tau$. It suffices to show that, for
each $T\ge 0$, the survival set
\[
A_T:=\{\vx\in X:\tau(\vx)>T\}
\]
is open in $X$. Fix $T\ge 0$ and $\vx_0\in A_T$. Then
$T<\tau(\vx_0)$, so by the definition of $\tau$ there exists
$R>\|\vx_0\|$ such that $T<\sigma_{K_R}(\vx_0)$. Then
\[
\sup_{0\le t\le T}\|\widehat{\vG}_{K_R}(\vx_0,t)\|<R.
\]
By continuity of the global flow $\widehat{\vG}_{K_R}$, uniformly over the
compact interval $[0,T]$, the same strict inequality holds for all initial
conditions $\vx$ in a neighborhood of $\vx_0$. Hence
$T<\sigma_{K_R}(\vx)\le \tau(\vx)$ for all such $\vx$, so this neighborhood is
contained in $A_T$. Therefore $A_T$ is open, proving that $\tau$ is lower
semicontinuous.

For continuity of $\vG$, fix $(\vx_0,t_0)\in D$. Choose
$T\in(t_0,\tau(\vx_0))$ and then choose $R>\|\vx_0\|$ with
$T<\sigma_{K_R}(\vx_0)$. By the preceding argument, for all $\vx$ near $\vx_0$
and all $t$ near $t_0$, the local flow agrees with the continuous global flow
$\widehat{\vG}_{K_R}$. Therefore $\vG$ is continuous on $D$.
\end{proof}

\lemLipschitz*
\begin{proof}[Proof of Lemma~\ref{lemma:Lipschitz}]
The global Lipschitz statement is Lemma~\ref{lemma:global-lip-flow}. For locally
Lipschitz $\vQ$, Lemmas~\ref{lemma:localized-global-systems}
and~\ref{lemma:localized-consistency} construct compatible global
localizations. Lemma~\ref{lemma:local-flow-construction} patches them into a
local flow. Lemma~\ref{lemma:maximal-local-flow} gives maximality and
uniqueness, and Lemma~\ref{lemma:lsc-flow-continuity} gives the lower
semicontinuous lifetime and continuity on the local-flow domain. These are
exactly the requirements in Assumption~\ref{ass:well-behaved-proj}.
\end{proof}

\subsection{Measurability of Dynamic Solvers}

\propMeasurable*
\begin{proof}\label{proof:measurable}
We prove measurability by decomposing $X$ according to the three possible
terminal behaviors of the local flow. Let $C_1$ be the set of initial
conditions whose trajectories have finite lifetime, let $C_2$ be the set of
initial conditions whose trajectories survive forever but do not converge, and
let $C_3$ be the set of initial conditions whose trajectories survive forever
and converge. On $C_1\cup C_2$ the dynamic solver returns $\dagger$; on $C_3$
it returns the limit of the trajectory. Thus, once we show that
$C_1,C_2,C_3$ are Borel and that the limit map on $C_3$ is Borel,
measurability follows from the representation
$g_{\vQ}(\vx)=\dagger$ on $C_1\cup C_2$ and
$g_{\vQ}(\vx)=L(\vx):=\lim_{t\to\infty}\vG(\vx,t)$ on $C_3$.

We now define these sets formally. The infinite-survival set is
\[
A_\infty:=\{\vx:\tau(\vx)=\infty\}=\bigcap_{n=1}^\infty A_n
\]
where $A_n:=\{\vx\in X:\tau(\vx)>n\}$. By lower semicontinuity of $\tau$, each
$A_n$ is open in $X$, so $A_\infty$ is Borel.

On $A_\infty$, the flow $\vG(\vx,t)$ is defined for every $t\ge 0$. This
restriction is important: outside $A_\infty$, the values $\vG(\vx,q)$ are not
defined for all large rational $q$, and replacing missing values by $\dagger$
would not give coordinate functions in $\R^d$. For $i\in\{1,\dots,d\}$, write
$G_i(\vx,t)$ for the $i$-th coordinate of $\vG(\vx,t)$ and define
extended-real maps on $A_\infty$ by
\begin{align*}
\ell_i^+(\vx)
&:= \inf_{T\in\N}\ \sup_{\substack{q\in\Q\\ q\ge T}} G_i(\vx,q),\\
\ell_i^-(\vx)
&:= \sup_{T\in\N}\ \inf_{\substack{q\in\Q\\ q\ge T}} G_i(\vx,q).
\end{align*}
For each rational $q\ge 0$, the map $\vx\mapsto G_i(\vx,q)$ is continuous on
$A_\infty$ because $A_\infty\subseteq A_q$ and $\vG(\cdot,q)$ is continuous on
$A_q$. Hence countable suprema and infima make $\ell_i^+$ and $\ell_i^-$
measurable on $A_\infty$. These maps are exactly
$\limsup_{t\to\infty}G_i(\vx,t)$ and
$\liminf_{t\to\infty}G_i(\vx,t)$, because
$t\mapsto G_i(\vx,t)$ is continuous on $[0,\infty)$ and suprema and infima over
real times can be taken over rational times.

The convergence set is
\[
\Xi:=A_\infty\cap
\bigcap_{i=1}^d
\{\vx:\ell_i^+(\vx)=\ell_i^-(\vx)\in\R\}.
\]
It is Borel in $A_\infty$, and therefore in $X$, because the defining
coordinate conditions are Borel. On $\Xi$, the coordinate limits exist. Define
their $i$-th coordinate by
$L_i(\vx):=\lim_{t\to\infty}G_i(\vx,t)=\ell_i^+(\vx)=\ell_i^-(\vx)$, and set
$L(\vx):=(L_1(\vx),\dots,L_d(\vx))$.
Each $L_i$ is measurable on $\Xi$, so $L:\Xi\to\R^d$ is Borel. Since $X$ is
closed and $\vG(\vx,t)\in X$, we have $L(\vx)\in X$ for every $\vx\in\Xi$.

Define the condition sets
\[
C_1:=X\setminus A_\infty,\qquad
C_2:=A_\infty\setminus\Xi,\qquad
C_3:=\Xi.
\]
By the definition of the dynamic solver, $g_{\vQ}|_{C_3}=L$.

For every Borel $B\subseteq X$, $g_{\vQ}^{-1}(B)=C_3\cap L^{-1}(B)$ is Borel
in $X$, while $g_{\vQ}^{-1}(\{\dagger\})=C_1\cup C_2=X\setminus C_3$ is Borel.
Hence
$g_{\vQ}:X\to X\sqcup\{\dagger\}$ is measurable.
\end{proof}

\subsection{Terminal Outcomes and Roots}

We now prove the containment underlying Proposition~\ref{prop:finite-rto-roots}.
Throughout this subsection, $\mR(\QTC)$ denotes the zero set of the projected
vector field $\QTC:X\to\R^d$.

The claim is intuitive: if a trajectory converges to an ordinary terminal
outcome, the vector field cannot keep pushing it in any feasible direction at the
limit. The proof uses the normal-cone formulation of projected dynamics and a
time-average argument. This avoids assuming that $\QTC$ is continuous at
boundary points.

\begin{lemma}\label{lemma:containment}
A dynamic solver $g_{\vQ}$ induced by $\vQ$ satisfies the following containment:
\[
g_{\vQ}(X) \subseteq \mR(\QTC) \cup \{\dagger\}.
\]
\end{lemma}
\begin{proof}\label{proof:containment}
Fix an initial condition $\vx_0\in X$ and write
$\vx(t):=\vG(\vx_0,t)$ for the corresponding projected trajectory. If the
system does not converge, then $g_{\vQ}(\vx_0)=\dagger$ and the claim is
immediate. Hence assume $g_{\vQ}(\vx_0)=\vx^*\in X$, so $\vx(t)\to \vx^*$.

\emph{Step 1: use the normal-cone inequality along the trajectory.}
For a.e. $t$, using the normal cone formulation, the projected dynamics satisfies:
\[
\vQ(\vx(t)) - \dot \vx(t)\in N_X(\vx(t))
\]
Throughout the proof, we consider times $t$ at which this relation holds. By definition of normal cone, for every $\vy\in X$,
\[
\begin{aligned}
&\langle \vQ(\vx(t))-\dot \vx(t),\,\vy-\vx(t)\rangle \le 0,\\
&\Rightarrow \langle \vQ(\vx(t)),\,\vy-\vx(t)\rangle\\
&\qquad\le \langle \dot \vx(t),\,\vy-\vx(t)\rangle.
\end{aligned}
\]
Notice
\[
\frac{d}{dt} \left(-\frac12 \|\vy-\vx(t)\|^2 \right)
=\langle \dot \vx(t),\vy-\vx(t)\rangle .
\]
Substituting this identity gives
\[
\langle \vQ(\vx(t)),\vy-\vx(t)\rangle
\le -\frac12\frac{d}{dt}\|\vy-\vx(t)\|^2
\]

\emph{Step 2: average over time.}
Integrating from $0$ to $T$ and dividing by $T$ yields
\[
\begin{aligned}
&\frac1T\int_0^T \langle \vQ(\vx(t)),\vy-\vx(t)\rangle\,dt\\
&\qquad\le \frac{\|\vy-\vx(0)\|^2-\|\vy-\vx(T)\|^2}{2T}.
\end{aligned}
\]
Taking $\limsup_{T\to\infty}$ on both sides, we get
\[
\begin{aligned}
&\limsup_{T\to\infty}
\frac1T\int_0^T \langle \vQ(\vx(t)),\vy-\vx(t)\rangle\,dt\\
&\qquad\le \limsup_{T\to\infty}
\frac{\|\vy-\vx(0)\|^2-\|\vy-\vx(T)\|^2}{2T}\\
&\qquad=0.
\end{aligned}
\]

Since $\vx(T)\to\vx^*$, the right-hand side is bounded divided by $T$ and hence
has limit $0$.

\emph{Step 3: pass to the limit.}
As \(\vQ\) and \(\vx\) are continuous, the scalar function
\(f(t):=\langle \vQ(\vx(t)),\vy-\vx(t)\rangle\) converges to
\(f(\infty):=\langle \vQ(\vx^*),\vy-\vx^*\rangle\) as \(t\to\infty\).
Moreover, since \(\vx(t)\to\vx^*\), the trajectory \(\vx(\cdot)\) is bounded,
and continuity of \(\vQ\) on the trajectory implies that \(f\) is bounded on
$[0,\infty)$. Hence the rescaled functions $s\mapsto f(Ts)$ on $[0,1]$ are
dominated by an integrable constant and converge pointwise a.e. to $f(\infty)$.
By the dominated convergence theorem (equivalently, the standard Ces\`aro averaging
lemma),
\begin{align*}
\lim_{T\to\infty}\frac{1}{T}\int_0^T f(t)\,dt
    &=\lim_{T\to\infty}\int_0^1 f(Ts)\,ds \\
    &=\int_0^1 f(\infty)\,ds \\
    &=f(\infty).
\end{align*}

Thus the time average converges to \(\langle \vQ(\vx^*),\vy-\vx^*\rangle\) as claimed.

Combining this with the previous inequality, we get
\begin{align*}
\langle \vQ(\vx^*),\vy-\vx^*\rangle \leq 0
\end{align*}
As this holds for every $\vy\in X$, we have
$\vQ(\vx^*) \in N_X(\vx^*)$. By definition of projected dynamics, this is
equivalent to $\QTC(\vx^*)=0$. Hence $\vx^* \in \mR(\QTC)$, proving the
containment.
\end{proof}
